\section{Estado do projeto}

\subsection{Decisões estabelecidas}

\begin{table}[ht]
\centering
\begin{tabular}{ll}
\toprule
\textbf{Parâmetro} & \textbf{Valor ou decisão} \\
\midrule
VLEN & 128 bits \\
ELEN & 32 bits \\
Elementos inteiros modelados & SEW de 8, 16 e 32 bits \\
LMUL inteiro modelado & 1, 2, 4 e 8 \\
Quantidade de registradores vetoriais físicos & 32 \\
Endianness da memória & Little-endian \\
Acesso vetorial à memória & Unit-stride, contíguo \\
Política de cauda & Tail-undisturbed \\
Política de máscara & Mask-undisturbed \\
Execução do modelo & Sequencial, funcional e sem temporização \\
\bottomrule
\end{tabular}
\caption{Parâmetros e decisões atuais do projeto.}
\end{table}

\subsection{Estado observado nos testes}

Na última execução registrada de \texttt{pytest -v}, foram coletados
38 testes: 31 passaram e 7 falharam. As falhas apontam para a mesma
causa imediata: o nome \texttt{SEW} não está definido em
\texttt{model/executor.py}, embora seja usado pelo handler de
\texttt{vsetvli}.

Portanto, os testes de estado, operações básicas, SEW/LMUL e geração
estática do programa de strip-mining passaram. Os testes que executam
programas pelo \texttt{Executor}, inclusive os testes end-to-end de
strip-mining, falharam ao processar a primeira instrução
\texttt{vsetvli}.

Este resultado não deve ser descrito como uma suíte totalmente aprovada.
A correção do import e a repetição da suíte são pendências antes de
considerar a especificação executável validada integralmente.